\documentclass[12pt,a4paper]{article}
\usepackage[utf8]{inputenc}
\usepackage[spanish]{babel}
\usepackage{amsmath,amssymb,amsfonts}
\usepackage{geometry}
\usepackage{hyperref}
\geometry{margin=2.5cm}

\title{\textbf{Constantes Físicas en Óptica}}
\author{Compendio de Constantes Fundamentales}
\date{\today}

\begin{document}

\maketitle

\section*{Introducción}
La óptica es la rama de la física que analiza el comportamiento, propagación y propiedades de la luz y la radiación electromagnética visible, ultravioleta e infrarroja. Las constantes en óptica determinan la velocidad de fase electromagnética en el vacío y las relaciones fundamentales de emisión de cuerpo negro.

\section{Velocidad de la Luz en el Vacío ($c$)}

\subsection*{1. Significado, Símbolo y Explicación}
Simbolizada con la letra $c$ (del latín \textit{celeritas}), es la velocidad máxima a la que la información, la materia y la radiación electromagnética pueden viajar en el universo. Es el nexo fundamental entre espacio y tiempo en la Teoría de la Relatividad Especial de Einstein y la relación fundamental de la energía ($E = mc^2$).

\subsection*{2. Valores Típicos en Ingeniería y Ciencia}
En cálculos habituales de óptica y física aplicada:
\begin{equation*}
c \approx 3 \times 10^8 \text{ m/s} \quad \text{o} \quad 2.998 \times 10^8 \text{ m/s}
\end{equation*}

\subsection*{3. Decimales Completos y Definición Exacta (SI 1983/2019)}
Desde la $17.^{\text{a}}$ Conferencia General de Pesas y Medidas (CGPM) en 1983, la velocidad de la luz se fijó como un valor exacto para definir la unidad fundamental de longitud (el metro):
\begin{equation*}
c = 299\,792\,458 \text{ m/s} \quad \text{(Exacto por definición)}
\end{equation*}
\textbf{Incertidumbre / Margen de Error:} Exacto sin incertidumbre ($u_r = 0$).

\vspace{0.8cm}
\hrule
\vspace{0.8cm}

\section{Primera Constante de Radiación ($c_1$ o $c_{1L}$)}

\subsection*{1. Significado, Símbolo y Explicación}
Representada como $c_1$ (o $c_{1L}$ para radiancia espectral), es la constante que aparece en la Ley de Planck de la radiación de cuerpo negro:
\begin{equation*}
c_1 = 2\pi h c^2 \quad \text{y} \quad c_{1L} = 2 h c^2
\end{equation*}
Determina la amplitud de la densidad espectral de emitiatividad radiante en función de la temperatura del emisor óptico.

\subsection*{2. Valores Típicos en Ingeniería y Ciencia}
\begin{equation*}
c_1 \approx 3.7418 \times 10^{-16} \text{ W}\cdot\text{m}^2 \quad \text{o} \quad c_{1L} \approx 1.191 \times 10^{-16} \text{ W}\cdot\text{m}^2\cdot\text{sr}^{-1}
\end{equation*}

\subsection*{3. Decimales Completos y Valor Exacto (SI 2019)}
Como es producto directo de la constante de Planck ($h$) y la velocidad de la luz ($c$), ambas exactas en el SI 2019:
\begin{equation*}
c_1 = 3.7417718521927573 \times 10^{-16} \text{ W}\cdot\text{m}^2
\end{equation*}
\begin{equation*}
c_{1L} = 1.1910429723971884 \times 10^{-16} \text{ W}\cdot\text{m}^2\cdot\text{sr}^{-1}
\end{equation*}
\textbf{Incertidumbre / Margen de Error:} Exacto por definición derivada ($u_r = 0$).

\end{document}
